\section{Algorithmic typing}

Algorithmic typing (\cref{fig:typing-algorithmic}). In both algorithmic
relations, $\InferType \E \T \Eff$ (inference) and $\CheckType \E \T \Eff$
(checking), $\Ctx$, $\Eff$ and $\CSet$ are outputs. In the inference relation
$\T$ is a further output but an input to the checking relation.

Correctness of splitting is checked via constraints accumulated in a multi
set-like structure $\CSet$. \[ \CSet \grmeq \CSEmpty \grmor \T \CEq \TU \grmor
\T \CSub \TU \grmor \CSet\CSJoin\CSet \]

We define the operator $\cdot \CtxJoin \cdot$ to join contexts according to the
direction $\Dir$. \[
  \Ctx[1] \CtxJoin \Ctx[2] = \begin{cases}
    \Ctx[1] \CtxPar \Ctx[2] & \text{if $\Dir \in \{ \DirUnord, \DirNone \}$} \\
    \Ctx[1] \CtxSeq \Ctx[2] & \text{if $\Dir = \DirLeft$} \\
    \Ctx[2] \CtxSeq \Ctx[1] & \text{if $\Dir = \DirRight$}
  \end{cases}
\]

The operator $\cdot|_{\cdot}$ restricts a context to a set of variables. \[
  \CtxRestrict*{X} = \begin{cases}
    \CtxRestrict*[1]{X} \CtxPar \CtxRestrict*[2]{X} & \text{if $\Ctx=\Ctx[1] \CtxPar \Ctx[2]$} \\
    \CtxRestrict*[1]{X} \CtxSeq \CtxRestrict*[2]{X} & \text{if $\Ctx=\Ctx[1] \CtxSeq \Ctx[2]$} \\
    \CtxVar \EVar \T & \text{if $\Ctx=\CtxVar \EVar T$ and $\EVar \not\in X$} \\
    \CtxEmpty & \text{else}
  \end{cases}
\]

\begin{figure}
  \hfill \fbox{$\InferType \E \T \Eff$}
   \quad \fbox{$\CheckType \E \T \Eff$}
   \begin{mathpar}
     \RuleAlgConst \and
     \RuleAlgVar \and
     \RuleAlgLetUnit \and
     \RuleAlgLetPair \and
     \RuleAlgApp \and
     \RuleAlgAbs \and
     \RuleAlgAbsRec \and
     \RuleAlgPair \and
     \RuleAlgCheck
%    \RuleTAppUnr \and
%    \RuleTAbsLin \and
%    \RuleTAppLin \and
%    \RuleTAbsLeft \and
%    \RuleTAppLeft \and
%    \RuleTAbsRight \and
%    \RuleTAppRight \and
%    \RuleTPairUnord \and
%    \RuleTPairOrd \and
%    \RuleTLet \and
%    \RuleTLetUnit \and
%    \RuleTLetUnord \and
%    \RuleTLetOrd \and
%    \RuleTInj \and
%    \RuleTCaseSum \and
%    % \RuleTFork \and
%    % \RuleTMatch \and
%    % \RuleTMatchLin \and
%    \RuleTWeakCE
%    % \RuleTWeakC
%    % \\
%    % \Big(\RuleTWeakC
%    % \and
%    % \RuleTWeakE\Big)
  \end{mathpar}
  \caption{Algorithmic typing}
  \label{fig:typing-algorithmic}
\end{figure}
